\section{证据约束下的 GPT-Live 架构重建}

公开证据可以确定 GPT-Live 在发声时持续监听、按细粒度时钟作出交互决策，并把复杂任务异步委派给 frontier model；但这些事实不足以推出其具体 tokenizer、fusion 位置或 action head。满足这些约束的最小重建是一种混合架构：media fast path 承载连续音频，stateful interaction core 维护双流因果状态，action channel 触发工具和委派，异步 cognition plane 负责慢速推理，而 authoritative state 负责转写、安全与 context compaction。

\subsection{哪些部分可以确定}

综合官方材料，可以确定以下系统性质：\FOne（1）语音模型在发声时继续监听，不以传统轮次检测器作为音频处理的推进条件；（2）模型每秒多次决定何时说、听、暂停、接受打断、调用工具或发起委派；（3）复杂任务可以异步委派，同时保持交互不中断；（4）媒体快速路径与应用逻辑由异步 RPC 隔离；（5）会话推理具有持久状态，并支持后台预填充和实例切换；（6）上下文压缩不会阻塞实时路径；（7）WebRTC/WARP、丢包处理与时钟漂移校正共同决定端到端体验；（8）系统容量以会话音频帧能否按时完成来衡量\cite{openai2026gptlive,openai2026continuous}。

这些性质共同排除了两种过度简化的解释。第一，GPT-Live 不太可能只是“更快的 ASR、LLM 与 TTS”，因为这类管线依赖离散轮次，无法自然处理系统发声期间到达的新输入。第二，它也不应被理解为“由一个超大模型承担全部任务”，因为官方明确说明系统会把复杂任务异步交给前沿模型处理。

\subsection{最可信但尚未证实的内部机制}

\Hyp 交互 core 很可能具有等价于双流时间格的内部状态：每个 chunk 编码新用户音频和近期 self-output，更新 causal state，再生成 assistant acoustic unit 与 interaction action。该表示可能像 Moshi 的 parallel stream、SyncLLM 的同步序列、LSLM 的 continuous encoder 加 fusion，也可能是未公开的新结构；公开信息不足以二选一。

\Hyp 为了在自身说话时理解用户，系统必须获得至少一种 self-speech reference：波形级 AEC 后音频、模型生成 token、播放时间戳或其组合。否则 residual echo 会和真实用户输入混淆。官方提到音频时钟与播放追赶，但未披露 echo reference 如何进入模型。

\Hyp “交互模型委派 frontier model”可以通过结构化事件实现：交互模型输出 delegate intent 和压缩后的 query，应用服务器把最终确认的 context 及 tool schema 送入已经预热的 session；返回结果写入 event stream，再由交互模型选择合适时机将其转化为自然口语。这个过程既不要求语音模型暴露内部思维，也不要求两个模型共享 KV cache。

\begin{figure}[H]
\centering
\resizebox{\textwidth}{!}{%
\begin{tikzpicture}[>=Latex,node distance=7mm and 9mm,font=\small,
box/.style={draw,rounded corners,align=center,minimum height=9mm},
f/.style={->,very thick,official},e/.style={->,thick,dashed,hypothesis}]
\node[box,fill=blue!7] (mic) {Audio capture\\AEC / jitter};
\node[box,fill=green!8,right=of mic] (enc) {Streaming user\\encoder/tokenizer};
\node[box,fill=green!8,right=of enc] (core) {Duplex interaction\\causal state};
\node[box,fill=green!8,right=of core] (dec) {Streaming speech\\decoder};
\node[box,fill=blue!7,right=of dec] (play) {Playout clock\\speaker};
\node[box,fill=orange!10,below=12mm of core] (action) {Interaction/action head\\pause, stop, delegate, tool};
\node[box,fill=orange!10,below=12mm of action] (frontier) {Asynchronous cognition\\reasoning + tools};
\node[box,fill=gray!10,left=12mm of frontier] (state) {Authoritative state\\transcript + safety};
\draw[f] (mic)--(enc); \draw[f] (enc)--(core); \draw[f] (core)--(dec); \draw[f] (dec)--(play);
\draw[e] (play) to[bend left=25] node[above,font=\scriptsize]{self-reference}(mic);
\draw[e,<->] (core)--(action); \draw[e,<->] (action)--(frontier); \draw[e,<->] (state)--(frontier); \draw[e] (state)|-(core);
\end{tikzpicture}%
}
\caption{可实现 GPT-Live 公开行为的最小架构重建。绿色语音内部结构与橙色动作头均属假设；官方确认的是持续交互、异步委派和系统边界。}
\label{fig:reconstruction}
\end{figure}

\subsection{“混合架构”比“纯端到端”更准确}

在感知--发声路径上，端到端学习可以保留韵律、重叠和非语言信号；在复杂任务路径上，模块化委派允许使用更强推理模型、工具权限和独立扩缩容；在 transport/state 层，确定性协议和显式状态机仍不可替代。因此 GPT-Live 更像“端到端交互模型嵌入模块化实时系统”，而不是对所有组件做单模型吞并。

这种拆分还能带来安全和成本方面的优势：小而快的交互模型可以高频运行，昂贵模型只在必要时启动；工具权限集中在可审计的认知平面；安全控制可以同时作用于实时输出和后台结果。代价则是跨层一致性、过期结果处理、委派策略和会话状态恢复会更加复杂。

\section{一个可复现的开源实现方案}

\subsection{最小可行架构}

建议以三个独立服务实现研究原型：

\begin{enumerate}
  \item \textbf{Realtime gateway：}WebRTC ingress/egress、Opus、AEC/jitter buffer、播放时间戳、连接与弱网遥测。媒体包进入有界 ring buffer，应用事件进入可靠 data channel。
  \item \textbf{Duplex inference worker：}包括 causal audio encoder、文本 LLM backbone、user-stream adapter、assistant audio decoder 和 action head。worker 保持固定的 session affinity，每 80 或 160 ms 执行一次 inference tick。
  \item \textbf{Agent 与 state service：}负责最终转写、conversation event log、tool runtime、reasoning model、context compaction、安全策略和 instance handoff 协调。
\end{enumerate}

原型可以从 Moshi 开源栈的语音主干出发\cite{defossez2024moshi}，增加独立的 action vocabulary 和 asynchronous tool bus；也可以采用 Freeze-Omni/LSLM 的路线，冻结文本 LLM，并接入 streaming encoder 和轻量 audio decoder，以降低训练成本\cite{wang2024freezeomni,ma2025lslm}。第一阶段不宜同时追求视频、超长 context 和多工具能力，应先验证真实 double-talk 条件下的 audio frame deadline 和打断行为。

\subsection{帧级调度伪代码}

\begin{algorithm}[H]
\caption{Deadline-aware duplex session loop}
\begin{algorithmic}[1]
\ForAll{active session $s$ in earliest-deadline order}
  \State $u_t\gets$ readAudioChunk$(s,\Delta)$
  \State $e_t\gets$ drainReadyEvents$(s,\mathrm{contextEpoch})$
  \State $(a_t,z_t,m_t)\gets$ duplexStep$(u_t,e_t,m_{t-1})$
  \If{$z_t$ requests delegation and no equivalent task is active}
    \State enqueueAsync$(taskId,\mathrm{contextEpoch},\mathrm{query})$
  \EndIf
  \If{safety permits $a_t$ and playout deadline remains}
    \State sendAudio$(s,a_t,\mathrm{timestamp})$
  \Else
    \State applyGracefulDegradation$(s)$
  \EndIf
\EndFor
\State run background prefill, compaction, and telemetry with residual budget
\end{algorithmic}
\end{algorithm}

调度循环需要避免三个隐蔽的阻塞点：tokenizer 或 audio codec 在 CPU 上串行执行；等待最终 ASR transcript 后才更新状态；在 media coroutine 中错误地同步等待 tool task。所有耗时操作都应通过 event 返回结果，media loop 只消费已经就绪的 event。

\subsection{分阶段实验计划}

\paragraph{阶段 A：离线可学习性。}
固定 7B--9B 文本 backbone 和 audio representation，在训练 token 数相同的条件下比较 channel fusion、cross-attention 和 gated fusion。训练集同时包含 half-duplex 对话、合成 overlap speech 和小规模真实双声道数据。验收标准包括：spoken QA 相对 half-duplex baseline 的退化不超过预设阈值；有效 barge-in 的 recall 提升且 false interruption rate 不恶化；存在 echo 时不发生 self-trigger。

\paragraph{阶段 B：单机实时性。}
在固定 GPU 上逐步增加并发 session，记录每个 tick 的 $Q/C$、DMR、P95/P99 jitter、KV cache 显存和 audio buffer underrun。比较 FIFO、continuous batching、earliest deadline first（EDF）以及预留 20\% 算力等策略。验收不以最高 token throughput 为准，而以每块 GPU 能稳定承载的 session 数为准。

\paragraph{阶段 C：异步认知。}
接入一个文本推理模型和三类工具：低延迟查询、长延迟搜索、多轮参数补全。对比不委派、阻塞式委派和异步委派三种模式，检查工具结果正确率、交互停顿、结果过期率、任务取消后的算力浪费，以及用户的主观等待感受。

\paragraph{阶段 D：长会话与故障。}
运行 30、60 和 120 分钟的长 session，主动注入 instance restart、network handoff、context compaction、transcript revision 和 tool-result race。验收标准包括：instance handoff 时不出现语音重复或缺失；context epoch 保持一致；用户关键事实得到保留；GPU 显存不随 session 时长无限增长。

\subsection{资源与延迟预算示例}

以 160 ms 为一个 inference tick 时，可以采用如下参考预算：uplink 与 jitter buffer 40 ms，streaming encoder 15 ms，interaction backbone 的 queueing 与 compute 45 ms，audio head 与 codec 20 ms，downlink、jitter buffer 和 playout 40 ms。这不是 GPT-Live 的公开参数，只是一个工程起点。预算必须使用 P95/P99 而不是平均值进行验证，并在跨地域、低端客户端和 5\% burst loss 条件下重新测试。

成本模型应按会话分钟分解为 interaction GPU、cognition GPU、音频带宽、tool 调用和 idle reservation。若委派率 $p_d$、平均 reasoning 时间 $T_r$、交互 core 每分钟成本 $C_f$，则粗略成本为
\begin{equation}
C_{min}=C_f+p_dT_rC_s+C_{net}+C_{tool}+C_{idle}.
\end{equation}
降低 $p_d$ 可以节省成本，但也可能损害任务正确率；因此，应结合委派准确率、召回率和最终任务成功率共同调优。

\section{局限性、开放问题与结论}

\subsection{仍然未知的关键细节}

本文不能回答 GPT-Live 的具体参数量、训练数据、codec、audio token rate、fusion 层位置、是否联合输出 action token、回声参考注入方式、推理并行策略和部署规模。这些信息从未被 OpenAI 官方公开。报告中的图~\ref{fig:reconstruction} 是满足公开行为的\emph{一种最小实现}，不是泄露或反向工程结论。

现有论文之间也很难公平横比。Moshi 的“约 200 ms”、Neural-FSM 的“500 ms 内响应”、GPT-4o 的“平均 320 ms”使用不同任务、网络、硬件和延迟起点。本文保留这些数字用于说明各自设计目标，不把它们排成统一排行榜。

\subsection{研究前沿}

未来最值得关注的方向包括：（1）在强语义整合和抗 overlap 干扰之间，自适应地选择 user-stream routing；（2）将 speech、language 和 action 放在统一 clock 上，同时避免高频 acoustic token 淹没稀疏 action；（3）利用真实双声道数据训练能够适应不同文化 turn-taking 习惯的模型；（4）为长 session compaction 建立事实保真和可逆审计机制；（5）联合优化 inference tick、GPU scheduling、network jitter 和 playout clock；（6）建立覆盖 stale result、task cancellation 和 provenance tracking 的 asynchronous delegation benchmark。

\subsection{结论}

GPT-Live 的核心并不是某个神秘的声学模块，而是将语音智能体重构为持续运行的实时系统。与轮次式全模态模型相比，它允许新输入在生成期间持续进入模型状态，并让交互策略与现实时间同步；与早期的单模型全双工方案相比，它进一步把复杂认知任务放到异步路径中处理，并通过有状态推理服务、后台上下文压缩与实例切换、低往返时延 WebRTC，将模型能力转化为稳定的产品体验。

因此，复现 GPT-Live 不应从“训练一个更大的语音 LLM”开始，而应同时建立四个闭环：\textbf{表示闭环}保证说和听能够并行，\textbf{控制闭环}学习停顿与打断，\textbf{认知闭环}安全地执行异步委派，\textbf{系统闭环}确保每个音频帧都能在播放时限之前到达。公开论文已经给出了这些组件的大部分候选机制；真正困难、也最有研究价值的，是让它们在同一个在线系统中稳定协同。
